\section{证据链：Upstream Quality 能否转化为 Downstream Quality}

本节从系统设计转向证据。若 sparse model 只是用更多总参数获得更低 pretraining loss，那么下游提升可能只是 upstream quality 的结果，而非参数形式本身。讲者因此提出控制问题：在相同 upstream quality 下，额外 sparse parameters 是否仍能独立改善 downstream task？

\lecturefigure{slide-27-upstream-quality-research-question.jpg}{研究问题：固定 upstream pretraining quality 后，总参数形式是否仍影响下游？}{Stanford Online 官方视频 00:43:58--00:44:17。}

\begin{knowledgebox}{背景概念：upstream 与 downstream}
Upstream task 是大规模预训练目标，例如 C4 上的 text-to-text denoising；downstream task 是 SuperGLUE、TriviaQA 等迁移评估。二者相关不意味着完全等价：同样 pretraining loss 的模型可能有不同表示、容量分配与 fine-tuning 行为。
\end{knowledgebox}

\subsection{SuperGLUE：相关性明显，但不是参数因果证明}

图中 dense 与 Switch 模型的 upstream quality 均与 SuperGLUE 提升相关。若两类点大体落在同一趋势附近，说明大量 downstream 收益可由更好的 pretraining quality 解释；若 sparse 点在同一 upstream 水平系统性更高，才支持参数形式有额外贡献。

\lecturefigure{slide-28-upstream-vs-superglue.jpg}{Dense 与 Switch 的 upstream quality 对 SuperGLUE quality 关系。}{Stanford Online 官方视频 00:44:18--00:45:19。}

\begin{knowledgebox}{读图：横向比较，不要只沿曲线看}
先在相近 upstream quality 的横坐标处比较 dense 与 sparse 点，再观察误差与模型规模。只看到两者都随 upstream 提升而上升，能说明相关性，却不能证明“更多稀疏参数在同质量下独立有效”。
\end{knowledgebox}

\begin{warningbox}{Correlation 不是机制识别}
模型族同时改变参数数、训练时间、正则化、路由与优化稳定性。散点趋势无法单独识别哪一项造成下游差异；需要 matched-quality、matched-compute 和多随机种子实验。
\end{warningbox}

Switch-C 将总参数推到约 1.6T，并在图中作为新的极端点。它的重要性是展示系统可训练性与 scaling 延续，而不是宣告所有任务都因“万亿”这个数字获得同比提升。极端点也可能对回归趋势产生很强杠杆作用，应与中等规模结果一起看。

\lecturefigure{slide-29-switch-c-trillion-parameter.jpg}{加入约 1.6T 参数的 Switch-C 后，观察 upstream 与 SuperGLUE 的扩展趋势。}{Stanford Online 官方视频 00:45:20--00:46:01。}

\teachervoice{讲者没有把 Switch-C 当作唯一卖点，而是用它测试趋势是否继续。工程上更重要的问题是：在相同设备时间、active compute 与服务约束下，它是否比更小的 sparse 或 dense 配置值得。}

\subsection{TriviaQA：knowledge-heavy 任务提供另一种观察}

TriviaQA 更偏向知识问答，因此适合检验“更多参数可能承载知识”的直觉。图中 upstream quality 与 TriviaQA 分数仍呈正相关，但一张相关图不足以把提升归因于存储容量；训练语料覆盖、tokenization、fine-tuning 与答案匹配规则都会影响结果。

\lecturefigure{slide-30-upstream-vs-triviaqa.jpg}{Upstream quality 与 TriviaQA 质量的关系。}{Stanford Online 官方视频 00:46:02--00:46:27。}

\begin{knowledgebox}{读图：这张图支持什么、不支持什么}
它支持“更好的预训练通常伴随更好的 TriviaQA”，并为 sparse capacity 值得继续研究提供证据；它不证明模型在内部以可定位方式存储事实，也不证明参数量比数据覆盖、检索或 active compute 更重要。
\end{knowledgebox}

\subsection{Base/Large fine-tuning 与 101 语言训练}

本节把证据从最大模型移回更常用规模与多语言分布。如果 sparse 方法只在最大模型上有效，它的工程价值会很窄。讲座展示 Base 与 Large 级别的 fine-tuning 结果，说明在更常见规模上也能获得收益；同时，多语言实验以 mT5 为 dense baseline，在 101 种语言上观察到整体优势。

\lecturefigure{slide-31-fine-tuning-base-and-large.jpg}{Base 与 Large 规模下的 fine-tuning 结果。}{Stanford Online 官方视频 00:46:28--00:47:07。}

\begin{knowledgebox}{读表：平均分之外还要找失败任务}
先比较同规模 dense 与 sparse 的平均质量和 FLOPs，再检查是否存在某些任务退化。平均提升可能掩盖对小数据、长序列或类别不平衡任务的敏感性；部署时应保留 per-task 结果，而不是只报 aggregate score。
\end{knowledgebox}

多语言环境天然具有异质输入，router 可能按语言、词形或主题形成路径分工；但这不意味着每种语言自动获得独立 expert。真正应检查的是高低资源语言的分布、tokenizer 覆盖、expert utilization 与是否出现某些语言被 overflow 系统性丢弃。

\lecturefigure{slide-32-multilingual-training.jpg}{以 mT5 为基线的 101 语言训练与质量趋势。}{Stanford Online 官方视频 00:47:08--00:47:41。}

\teachervoice{讲者把“所有 101 种语言都有改善”作为重要结果，但工程复现仍应分语言报告。宏平均改善不等于低资源语言、公平性和路由覆盖都已解决。}

\begin{warningbox}{语言分工可能制造新的长尾}
若 router 把高资源语言集中到稳定 experts，低资源语言可能得到更少训练或更高 dropping。应按语言统计 expert entropy、overflow、token 数和 downstream quality，避免总体吞吐掩盖局部退化。
\end{warningbox}

\subsection{本章小结}

Upstream quality 是 downstream quality 的重要解释变量，但相关图不自动证明额外稀疏参数的独立因果作用。Base/Large 与多语言结果扩大了适用范围，同时要求更细粒度的任务和语言审计。

